% 第 1 章 项目背景与问题分析 (目标 3 页)
% 对应官方要求: 问题分析
% 要点: 威胁三阶段演进 → 现有方案三类不足 → 关键缺口总结
\section{项目背景与问题分析}

\subsection{智能体安全的威胁场景演进}

大模型智能体进入政企实际业务后,工具调用能力成为标配。工具集合覆盖邮件外发、文件读写、数据库查询、日历创建、定时任务、系统命令执行等高权限动作。围绕这些动作的安全威胁,经历三阶段演变。

第一阶段为显式注入。攻击者在用户消息或上传文档中嵌入"忽略上述规则"等字面指令,或用方括号标记注入载体。结构化关键词匹配对该类攻击有较好的识别能力。

第二阶段为去标记化业务话术诱导。攻击者将诱导意图嵌入正常公文格式、会议纪要、邮件正文中,完全消除方括号、感叹号等明显攻击信号。LLM 在阅读上下文后将诱导意图误判为业务推进,主动调用外发或写入工具。本项目在自有评测集上测得,基线阶段 66 条去标记化业务话术样本中 57 条成功绕过 LLM 自身安全过滤,绕过率 86.4\%。

第三阶段为过度代理,也是政企场景中最具迷惑性的威胁形式。LLM 在完成用户原始任务后,自主决定顺手做额外的外发与写入。例如用户只要求查看差旅费执行情况,LLM 自主决定起草报告并发送给局领导。该类威胁在常见基准中往往被归为无攻击类别,实际造成真实数据外发。OWASP LLM Top 10 将其列为 LLM08 Excessive Agency 独立风险项。

% TODO: 图 1 威胁演进示意图(三阶段时间线+每阶段代表样本片段)
\placeholderfig{威胁三阶段演进示意(显式注入 → 业务话术诱导 → 过度代理)}{智能体工具调用威胁的三阶段演进}

\subsection{现有方案的不足}

针对上述三类威胁,公开方案与本项目基线阶段的实测暴露三类不足。

第一,关键词层对去标记化场景失效。本项目基线阶段使用关键词加短语匹配,典型词表覆盖"忽略规则""立即执行""伪造签字""删除审计"等。在 66 条去标记化样本上,该层仅命中 5 条。

第二,LLM 自身安全机制在去标记化场景下漏过比例高。66 条样本中 9 条被 LLM 自身拒答,3 条未触发攻击工具,57 条成功绕过。

第三,LLM 评审类方案在工程实现上存在开放失败缺陷,即上游调用异常或模型无响应时,系统倾向放行。该模式在通用场景可接受,在政务数据敏感场景不可接受。本项目基线阶段实测,因上游接口路径配置错误,全部 LLM 评审调用返回异常,全部异常均被放行,放行率与无评审状态一致。

\subsection{关键缺口总结}

现有方案在去标记化业务话术诱导、过度代理、评审工程安全三处存在明确缺口。三者协同缺失,导致真实业务话术下的攻击绕过率达 86.4\%。政企场景另有其特殊约束:数据敏感性要求异常时阻断而非放行,业务连续性要求正常操作零误伤,审计合规要求每一次判定留痕可复现。任何单点方案无法同时满足以上约束,需要体系化设计。
